%==============================================================================
% appendices/B_independent_proofs.tex
%==============================================================================
\section{Proofs for the Independent Baseline}
\label{appx:independent-proofs}

This appendix collects the full proofs of the results stated in
\cref{sec:independent-baseline}. Each proof environment is labelled by the
result it establishes.

\begin{proof}[Proof of \cref{thm:catastrophe-exact}]
Let \(A_i(x)=\{X_i>x\}\). By finite inclusion--exclusion,
\[
  \sum_i\Pb(A_i)-\sum_{i<j}\Pb(A_i\cap A_j)
  \le \Pb\!\left(\bigcup_iA_i\right)
  \le \sum_i\Pb(A_i).
\]
Under \cref{ass:ind-pte} the variables are independent, so
\[
  \Pb(A_i\cap A_j)=\Pb(X_i>x)\Pb(X_j>x)
  \sim c_ic_j\Gb(x)^2=o(\Gb(x)).
\]
There are finitely many pairs, so \(\sum_{i<j}\Pb(A_i\cap A_j)=o(\Gb(x))\).
Higher-order intersections are bounded above by pair intersections and so are
also \(o(\Gb(x))\). Combined with the marginal sum
\(\sum_i\Pb(X_i>x)\sim C\Gb(x)\), this yields
\(\Pb(M_N>x)\sim C\Gb(x)\).
\end{proof}

\begin{proof}[Proof of \cref{thm:sum-equivalence}]
Let \(h=h(x)\uparrow\infty\) with \(h=o(x)\) as in \cref{ass:ind-sbj}.
Upper bound. Decompose
\[
  \{S_N>x\}\subseteq
  \bigcup_{i=1}^N\{X_i>x-h\}
  \;\cup\;
  \left\{S_N>x,\ \max_i X_i\le x-h\right\}.
\]
The first term satisfies
\[
  \Pb\!\left(\bigcup_i\{X_i>x-h\}\right)
  \sim \sum_i\Pb(X_i>x-h)
  \sim \sum_i\Pb(X_i>x),
\]
by inclusion--exclusion and the local invariance
\(\Pb(X_i>x-h)\sim\Pb(X_i>x)\) of \cref{ass:ind-sbj}.
The second term is excluded by the truncation conditions in
\cref{ass:ind-sbj}: the joint event that two or more truncated components
exceed order \(h\) is \(o(\Gb(x))\), and the residual non-large contribution is
\(o(\Gb(x))\) by the same truncation.

Lower bound. With the same \(h\),
\[
  \Pb(S_N>x)
  \ge \sum_i\Pb(X_i>x+h,\ S_{-i}>-h) - o(\Gb(x)),
\]
where the overlap correction is \(o(\Gb(x))\) by the pair-exceedance estimate
of \cref{thm:catastrophe-exact}. Independence gives
\[
  \Pb(X_i>x+h,S_{-i}>-h)
  =\Pb(X_i>x+h)\,\Pb(S_{-i}>-h)
  \sim \Pb(X_i>x),
\]
since \(\Pb(S_{-i}\le -h)=o(1)\) and \(\Pb(X_i>x+h)\sim\Pb(X_i>x)\). Summing
yields the matching lower bound.
\end{proof}

\begin{proof}[Proof of \cref{thm:second-order}]
Set \(r(x)=\Gb(x)(|A(x)|\vee x^{-1})\). The events
\(\{i\ \text{selected}\}=\{R_i(X)=1\}\) partition \(\{S_N>x\}\) under the fixed
tie-breaking rule of \cref{ass:tie}. On the \(i\)th stratum
\[
  \Pb(S_N>x, R_i(X)=1)
  =\E\!\left[\Pb(X_i>T_i(x)\mid X_{-i})\,\1\{T_i(x)>M_{-i}\}\right].
\]
By \cref{ass:ind-sbj} the contribution to the right-hand side from the event
\(\{T_i(x)\le M_{-i}\}\) is \(o(r(x))\), and the same truncation gives
negligible mass on \(\{|S_{-i}|>h(x)\}\). It therefore suffices to expand
\(\E[\Pb(X_i>x-S_{-i})]\) on \(\{|S_{-i}|\le h(x)\}\) and then take
expectations.

By \cref{ass:sote}, locally uniformly on \(|S_{-i}|\le h(x)\),
\[
  \Pb(X_i>x-S_{-i})
  =\Gb(x)\!\left[c_i+c_id_iA(x)+\frac{\alpha c_iS_{-i}}{x}
  + o\!\left(|A(x)|+x^{-1}\right)\right].
\]
Taking expectations,
\[
  \E\Pb(X_i>x-S_{-i})
  = c_i\Gb(x)+c_id_i\Gb(x)A(x)
  + \frac{\alpha c_i\Gb(x)}{x}\,\E S_{-i}
  + o(r(x)),
\]
and \(\E S_{-i}=\mu_{-i}=\sum_{j\ne i}\E X_j\) by \cref{ass:sote}. Summing over
active \(i\) and discarding inactive contributions (which are lower order by
\cref{ass:sote}) yields the expansion in \cref{thm:second-order}.

\emph{Consistency check at \(N=1\).} If \(N=1\), then \(S_{-1}=0\) and
\(\mu_{-1}=0\); the mean-shift term vanishes, the marginal second-order term
reduces to \(c_1d_1\Gb(x)A(x)\), and the formula collapses to the standard
single-variable second-order expansion.
\end{proof}

\begin{proof}[Independent CdMC: unbiasedness and BRE bound]
\emph{Unbiasedness.} For continuous variables the smallest-index tie-breaking
rule of \cref{ass:tie} gives the disjoint union
\[
  \{S_N>x\}=\bigsqcup_{i=1}^N\{S_N>x,\ R_i(X)=1\}.
\]
Conditional on \(X_{-i}\), the \(i\)th stratum equals
\(\{X_i>T_i(x)\}\cap\{T_i(x)>M_{-i}\}\), and independence yields
\(\Pb(X_i>T_i(x)\mid X_{-i})=\Fb_i(T_i(x))\). Taking expectations and summing
over \(i\) gives \(\E Z^{\mathrm{ind}}(x)=\Pb(S_N>x)\).

\emph{Scale-invariant BRE bound.} We claim \(T_i(x)\ge x/N\) almost surely:
suppose \(M_{-i}<x/N\); then if also \(x-S_{-i}<x/N\), one has
\(S_{-i}>(N-1)x/N\). Since \(S_{-i}\) is a sum of \(N-1\) terms each bounded by
\(M_{-i}<x/N\), this gives \(S_{-i}\le (N-1)x/N\), a contradiction. Hence
\(T_i(x)\ge x/N\), and therefore
\[
  0\le Z^{\mathrm{ind}}(x)\le B(x):=\sum_{i=1}^N\Fb_i(x/N).
\]
Since \(\E Z^{\mathrm{ind}}(x)=\mu(x)\),
\[
  \Var Z^{\mathrm{ind}}(x)
  \le \E\!\left[Z^{\mathrm{ind}}(x)^2\right]-\mu(x)^2
  \le B(x)\mu(x)-\mu(x)^2.
\]
Regular variation gives \(\Fb_i(x/N)/\Fb_i(x)\to N^\alpha\) and the sum
equivalence (\cref{thm:sum-equivalence}) gives \(\sum_i\Fb_i(x)\sim C\Gb(x)\),
so \(B(x)/\mu(x)\to N^\alpha\). Consequently
\[
  \limsup_{x\to\infty}\frac{\Var Z^{\mathrm{ind}}(x)}{\mu(x)^2}
  \le N^\alpha-1.
\]
The bound is scale-invariant: replacing \(X\) by \(\lambda X\) leaves \(N\),
\(\alpha\), and the ratio unchanged.
\end{proof}

\subsection{Verification notes}

\paragraph{Lomax second-order constant.}
For the Lomax/Pareto-type benchmark $\Fb_i(t)=(1+t/\beta_i)^{-\alpha}$ on
$t>0$, $\E X_i=\beta_i/(\alpha-1)$ when $\alpha>1$. A direct substitution
verifies the cross-term in the second-order expansion of
$\Pb(X_i>x-S_{-i})$: expanding to second order in $S_{-i}/x$ gives
\[
  (1+(x-S_{-i})/\beta_i)^{-\alpha}
  = (x/\beta_i)^{-\alpha}\left(1-\beta_i/x\right)^{-\alpha}
    \left(1+S_{-i}/(x-\beta_i)\right)^{\alpha},
\]
and expansion to first order in $S_{-i}/x$ and to leading order in the
marginal-tail auxiliary $A(x)=\beta_i/x$ recovers
$c_i d_i A(x)+\alpha c_i S_{-i}/x$ with the same sign convention used in
\cref{thm:second-order}. This is the regression check used by the
implementation pipeline.

\paragraph{Scale-invariance check.}
The BRE bound $N^\alpha-1$ is invariant under $X_i\to\lambda X_i$ because both
$T_i(x)$ and $x$ scale by $\lambda$, leaving $\Fb_i(T_i(\lambda x))=\Fb_i(T_i(x))$
when $\Fb_i$ is correctly rescaled. The bound therefore depends only on $N$ and
$\alpha$, not on the marginal scale constants $c_i$.
